% Generated by report_gen.py (hwde v1) - do not hand-edit.
\documentclass[11pt,a4paper]{article}
\usepackage[margin=2.2cm]{geometry}
\usepackage{graphicx}
\usepackage{booktabs}
\usepackage{longtable}
\usepackage{pdfpages}
\makeatletter
\@ifundefined{KV@Gin@artifact}{\define@key{Gin}{artifact}[]{}}{}
\makeatother
\usepackage{xcolor}
\usepackage[hidelinks]{hyperref}
\setcounter{tocdepth}{1}
\makeatletter
\def\ps@hwde{\let\@mkboth\@gobbletwo\def\@oddhead{}\def\@evenhead{}\def\@oddfoot{\small PCB-0017-A -- stereo-class-d-amp\hfil\thepage}\let\@evenfoot\@oddfoot}
\makeatother
\pagestyle{hwde}
\begin{document}
\sloppy
\begin{center}
{\LARGE\bfseries stereo-class-d-amp --- Design Document}\\[6pt]
{\Large\bfseries PCB-0017-A}\\[6pt]
{\large hwde v1 pipeline}\\[2pt]
generated 2026-09-27 14:37:20
\end{center}
\tableofcontents

\section{Board and Run Metadata}
{\small
\begin{longtable}{lp{11cm}}
\toprule
\textbf{Field} & \textbf{Value} \\
\midrule
\endhead
board & stereo-class-d-amp \\
part number & PCB-0017-A \\
workspace & \texttt{boards/PCB-0017-A\_stereo-class-d-amp} \\
phase & P7 \\
gate status & 4/5 gates pass (not passing: verify) \\
state created & 2026-09-27T00:25:28 \\
state updated & 2026-09-27T08:57:51 \\
\bottomrule
\end{longtable}
}
\section{Overview}
\subsection*{PCB-0017-A stereo-class-d-amp ? brief}
Owner brief: rebuild and compare against a published 44 x 34 mm reference design.

\begin{itemize}
\item TPA3116D2 stereo class-D amplifier on a single 12 V supply, stereo input and output, SMD parts only,
2-layer board as small as possible. Target: match or beat the reference design's 44 x 34 mm.
\item Correctness and the TPA3116D2's thermal pad / heat-sinking come before size; say where we deviate.
\item The reference design's reported parts: TPA3116D2; 2.2 uF X7R input coupling; 100 uF + 100 nF + 1 nF supply bypass;
10 uH LC output filters per channel with damping branches; 5.6k/100k config and pull-up resistors;
wire pads J1-J4 and TP1. About 40 SMD parts, 45 footprints.
\item Source: https://jlcpcb.com/blog/gpt-6-astra-pcb-design-in-kicad
\item Stop at a DRC/ERC-clean fab package. Do NOT order. End with COMPARISON.md: our size, layers,
part count and DRC result against the reference design's.
\end{itemize}
\section{Requirements}
\subsection*{PCB-0017-A stereo-class-d-amp - requirements}
Source: \texttt{brief/brief.md} (owner brief). No mode token: normal design, full scope. Run is unattended; every open question in section 9 carries a default the orchestrator may take as the owner's delegate.

\subsubsection*{1. Function}
A stereo class-D audio amplifier board around the TI TPA3116D2, fed from a single 12 V supply, with stereo line-level input and stereo speaker output on wire pads. It reproduces a published 44 x 34 mm reference design (JLCPCB blog post, link in the brief) so we can compare our result against the reference design's: a 2-layer board, SMD parts only, as small as possible, with correctness and the TPA3116D2's thermal pad and heat-sinking taking priority over size. Work stops at a DRC/ERC-clean fab package (no order) plus a \texttt{COMPARISON.md} giving our size, layer count, part count and DRC result against the reference design's reported 44 x 34 mm, 2 layers, about 40 SMD parts and 45 footprints. Where we deviate from the reference design we say where and why.

\subsubsection*{2. Interfaces}
\begin{itemize}
\item Audio input, stereo, line level, single-ended: two channels plus ground, AC-coupled into the
TPA3116D2 through 2.2 uF X7R capacitors (the reference design's value). Connector: wire pads (the reference design's J1-J4
are wire pads, not connectors). Input impedance and level are the amplifier's defaults
(gain set by the 5.6k / 100k resistors the reference design reports). Which of J1-J4 is input vs output vs
power is not stated - see section 9 Q1.
\item Audio output, stereo, bridge-tied load (BTL) speaker outputs, one 10 uH LC output filter per
channel with a damping branch (the reference design's parts). Wire pads. Speaker impedance not stated - Q2.
\item Power input, 12 V DC, wire pads. Q3 covers a screw terminal or barrel jack instead.
\item TP1: one test point (the reference design's list). Purpose not stated; ASSUMED: a ground or supply probe pad.
\item Configuration: gain / mode set by fixed resistors (the reference design's 5.6k / 100k). No switches, no LEDs,
no buttons stated. Fault / shutdown / mute pins: pull-up resistors only (the reference design's list), no user
control - Q4 asks whether a mute or fault indicator is wanted.
\item No USB, no RF, no digital interface.
\end{itemize}
\subsubsection*{3. Power}
\begin{itemize}
\item Input: single 12 V DC supply (brief). Range not stated; ASSUMED 12 V nominal, tolerate roughly
10.8 to 13.2 V (a regulated 12 V adapter). The TPA3116D2 accepts a wide range, but the input
connector, bulk capacitor voltage rating and reverse-polarity choice hinge on the range - Q5.
\item Rails: PVCC = 12 V direct to the amplifier. The TPA3116D2 generates its own internal analog and
gate-drive rails; no external regulator is expected. GUESS: no other rail on the board.
\item Supply bypass per the reference design: 100 uF bulk + 100 nF + 1 nF at each PVCC group (the reference design reports one set;
the datasheet-driven count per pin group is P1's to settle).
\item Supply current GUESS (not researched, from the part's class): at 12 V the TPA3116D2 delivers on
the order of 15 W per channel into 4 ohm or about 8 W per channel into 8 ohm at the onset of
clipping. Both channels driven at full output into 4 ohm draws about 2.5 to 3 A from 12 V;
typical music program is well under 1 A. Whether this is treated as \textgreater{}3 A is a section 8 flag - Q6.
\item No battery, no charging.
\item Reverse-polarity or overvoltage protection at the input: not stated - Q7.
\end{itemize}
\subsubsection*{4. Environment}
Not stated. ASSUMED: indoor bench / hobby use, 0 to 40 C ambient, no enclosure specified, no ingress rating, no vibration requirement. The TPA3116D2 dissipates a few watts at full output, so thermal design assumes free air and relies on the PCB copper (thermal pad vias to a bottom copper pour) rather than an attached heatsink - Q8 confirms no clip-on heatsink is required.

\subsubsection*{5. Size \& mounting}
\begin{itemize}
\item Outline: as small as possible; target to match or beat the reference design's 44 x 34 mm on 2 layers. This is a
soft target, no HARD cap: correctness and the thermal pad / heat-sinking area come first, and the
brief asks us to state where we deviate. If the thermal copper needs more area than 44 x 34 mm
allows, the board grows and COMPARISON.md says so.
\item Layers: 2 (brief). Treated as a firm design choice, not a physical cap; P1 may argue for more only
with a stated reason and the deviation recorded.
\item Mounting holes: not stated. The reference design's board is not reported to have any. ASSUMED none (bare module
with wire pads), see Q9.
\item Height limit: none stated. The 100 uF bulk capacitor is the tallest part; SMD only per the brief.
\item SMD parts only (brief). Wire pads J1-J4 and TP1 are copper pads, no through-hole connectors.
\end{itemize}
\subsubsection*{6. Quantity \& budget}
Not stated. Nothing is ordered in this run (brief: stop at the fab package). ASSUMED for costing and DFM: 5 boards, hobby budget, no unit-cost target - Q10.

\subsubsection*{7. Assembly}
Not stated. ASSUMED: JLC PCBA, all parts on the top side (single-sided assembly) where the layout allows; the bottom side kept for the thermal copper pour and vias only. A double-sided assembly is allowed if it is the only way to hit the size target, and COMPARISON.md then says so - Q11. Parts should be JLC-stocked (basic / preferred where a choice exists) so the package is orderable as-is even though we do not order.

\subsubsection*{8. Compliance/safety flags}
\begin{itemize}
\item Mains voltage: no. The board takes 12 V DC from an external supply.
\item Batteries: no.
\item Motors: no.
\item \textgreater{}30 V: no (12 V supply; the TPA3116D2's own switching stays under 30 V).
\item High current (\textgreater{}3 A): borderline - flagged. The 12 V input and the BTL speaker outputs carry
amplifier current. At 12 V into 4 ohm, both channels at clipping, the supply draw is estimated
near 3 A (section 3 GUESS); the output filter inductors and traces carry a few amps peak per
channel. Q6 asks how to treat it. Not a guess to proceed on: the trace widths, inductor
saturation rating and pad sizing depend on the answer.
\item RF transmit: no. Class-D switching (about 400 kHz to 1.2 MHz) radiates; the LC output filter
and layout handle EMI but no intentional transmitter is present.
\item Speaker output is a bridge-tied load: neither output terminal is ground. Note for the wire pad
labelling and for anyone wiring it; not a safety hazard at these voltages.
\end{itemize}
\subsubsection*{9. Open questions}
1. Which of the wire pads J1-J4 carry what? Default: J1 = 12 V power in (V+ and GND), J2 = stereo line input (L, R, GND), J3 = left speaker out (+/-), J4 = right speaker out (+/-). Each "pad" is a small group of copper pads. 2. Speaker impedance the amplifier should be designed for? Default: 4 ohm minimum, 8 ohm typical (the LC output filter values are tuned for this; 10 uH per the reference design suits both). 3. Should the power input stay as bare wire pads like the reference design's, or be a real connector (screw terminal or barrel jack)? Default: wire pads, to match the reference design and the SMD-only rule. 4. Do you want a mute/standby control or a fault indicator LED? Default: no - fixed pull-ups so the amp is always on, matching the reference design's parts list. 5. Supply voltage range the board must accept? Default: 12 V nominal, 10.8 to 13.2 V (a regulated 12 V adapter); capacitors rated 25 V or better. 6. Current handling (safety flag, section 8): design the 12 V input path and speaker outputs for how much continuous current? Default: 3 A continuous on the 12 V input, 3 A per speaker output, which covers both channels driven fully into 4 ohm at 12 V. Answering "yes to the default" accepts that the board is sized for up to 3 A; a higher figure widens traces and grows the board. 7. Reverse-polarity protection on the 12 V input? Default: none (matches the reference design's about-40-part count); the user must wire polarity correctly. 8. Is PCB copper alone acceptable for heat-sinking (thermal pad vias to a bottom copper pour), with no attached heatsink? Default: yes - the TPA3116D2 at 12 V into 8 ohm runs cool on copper alone; into 4 ohm continuously it may need more copper area than 44 x 34 mm, and we would then grow the board and record the deviation. 9. Mounting holes? Default: none (bare module, like the reference design's). 10. Quantity and budget for the DFM / cost sanity check? Default: 5 boards, no unit-cost target, JLC-stocked parts. 11. Assembly: JLC PCBA, top side only? Default: yes, single-sided top; allow bottom-side parts only if that is the sole way to meet the size target, and say so in COMPARISON.md. 12. Input level and gain: line level (about 1 Vrms) with the gain the reference design's 5.6k/100k resistors set? Default: yes, keep the reference design's gain setting so the comparison is like-for-like.

\paragraph*{Answers to section 9 (orchestrator as owner's delegate, 2026-09-27; recorded as state.json decisions)}
All twelve defaults above are ACCEPTED as written. Q6 (current) is PROVISIONAL: 3 A continuous on the 12 V input and 3 A per speaker output; re-confirmed at H1 against the P1 power research.

\section{Architecture}
\subsection*{Decisions of record (state.json)}
{\small
\begin{longtable}{lp{5.2cm}p{7.2cm}}
\toprule
\textbf{Phase} & \textbf{What} & \textbf{Why} \\
\midrule
\endhead
P0 & Q1 pads: J1=12V in (V+,GND), J2=stereo line in (L,R,GND), J3=left spk +/-, J4=right spk +/-; TP1 per the reference design & delegate answer (owner said decide yourself); matches the reference design's J1-J4 \\
P0 & Q2 speaker: 4 ohm min, 8 ohm typ & delegate; 10uH filter suits both \\
P0 & Q3 power in: bare wire pads & delegate; like-for-like with the reference design \\
P0 & Q4 no mute/fault LED; fixed pull-ups, always on & delegate; the reference design has none \\
P0 & Q5 supply 12V nom, 10.8-13.2V, caps \textgreater{}=25V & delegate default \\
P0 & Q6 PROVISIONAL (safety, \textgreater{}3A borderline): 3A continuous on 12V input, 3A per speaker output & delegate; TPA3116D2 \textasciitilde{}2x15W into 4ohm at 12V -\textgreater{} \textasciitilde{}2.8A supply; unattended run, reconfirm at H1 \\
P0 & Q7 no reverse-polarity protection & delegate; matches the reference design part count; noted as a deviation risk in COMPARISON \\
P0 & Q8 PCB copper heat-sinking only (PowerPAD thermal vias to bottom pour), grow board if needed and record & owner: thermal first, then size \\
P0 & Q9 no mounting holes & delegate; bare module like the reference design \\
P0 & Q10 qty 5, no unit-cost target, JLC stocked parts & delegate default \\
P0 & Q11 JLC PCBA top side only unless size forces otherwise & delegate default \\
P0 & Q12 line-level in, keep the reference design's 5.6k/100k gain/config & delegate; like-for-like \\
P1 & DEVIATION: amp IC TPA3116D2DAD -\textgreater{} TPA3118D2DAP (same SLOS708 family, same 32-pin HTSSOP pinout, PowerPAD DOWN) & TPA3116D2 exists only as DAD (pad UP, TI SLOS708); its 14 C/W needs a top heatsink, which the reference design's own write-up admits is not modeled. Brief: thermal pad/heat-sinking first, SMD only. At 12 V output is rail-limited so no audio power lost; pad-down lets 2-layer copper + vias sink it (\textasciitilde{}22-30 C/W). \\
P1 & Q6 revised: +12V input and its GND return sized for 4 A; 3 A per speaker output; inductor Isat \textgreater{}= 4 A & power research: 3.6 A at 10\% THD into 4 ohm, 4.0 A at 13.2 V \\
P1 & Sustained full-power 4 ohm sine at 40 C is NOT a requirement (music + 8 ohm full power are); keep the size target & power research: that case needs \textasciitilde{}50x50 mm 2 oz or PLIMIT; thermal shutdown is the chip's own protection; FAULTZ-\textgreater{}SDZ gives auto-recovery \\
P1 & Prefer 2 oz outer copper if the stackup table offers it for 2L & halves trace widths and improves spreading; 0.65 mm pitch is fine at 2 oz floors \\
P1 & PLIMIT tied to GVDD (no limit); FAULTZ tied to SDZ with 100k pull-up for self-clearing thermal fault; gain 20 dB via 5.6k master & TI 12 V design; refdesign research resolves the reference design's 5.6k/100k as gain leg + pull-ups \\
P1 & No reverse-polarity/TVS part (Q7 stands) & like-for-like with the reference design; risk (reversed supply, hot-plug ring) written in COMPARISON.md \\
P1 & 10 uF PVCC decoupling in 0805 (Basic CL21A106KAYNNNE), not 1206 & Basic part, smaller footprint helps size; 25 V rating kept \\
P1 & No switch-node EMI snubber (10R+330pF); keep only the 3.3R+10nF damping branch per channel output & like-for-like with the reference design and the EVM damping network; EMI not a stated requirement \\
P1 & Parts leads: TPA3118D2DAPR C46497; SLO0630H100MTT 10uH C207842 (Isat 5.5A); 680nF out filter; bulk MA25V100M6x6 polymer + RVT 100uF 35V electrolytic & P1 component-scout shortlist \\
P2 & P2 research tiers: class-d-amp + lc-output-filter researchers on fable/high (tier); dc-input + audio-line-in researchers and second readers on sonnet/high (explicit downgrade) & B1/B2 are low-risk textbook blocks already covered by P1 power.md/refdesign; owner asked for quality without burning money \\
P2 & Start P3 part-sourcer in parallel with P2 research, before H1 & part choice doesn't turn on the checklist research; a value the research changes is amended in parts.json before P4 \\
P3 & C11/C12 10 uF PVCC: accept X5R CL21A106KAYNNNE (C15850) though sheets.md said X7R & caps sit below 85 C (X5R limit) at \textless{}=1.5 W design dissipation; DC-bias loss at 12 V is covered by the 100 uF bulk per pair; Basic part \\
P3 & R4-R7 3.3R 0603 100 mW (RC0603FR-073R3L) accepted & branch dissipation \textasciitilde{}1 mW: 10 nF is \textasciitilde{}800 ohm at 20 kHz and \textasciitilde{}40 ohm at 400 kHz against \textless{}0.3 V residual ripple \\
P2 & R4-R7+C23-C26 (3.3R+10nF) stay on the speaker side of the LC filter as TI's EMI C-RC snubber, not at the switch node & research: TI labels it an EMI snubber (corner \textasciitilde{}4.8 MHz), not audio damping; at the switch node 10 nF x 12 V\textasciicircum{}2 x 400 kHz = 0.58 W would overload a 0603. Filter Q 0.52-1.04 over 4-8 ohm is acceptable per SLAA701B \\
P3 & designing under coverage gap: block:B1 dc-input classes power-loop, esd, sourcing, constraints-emission have no record at the approved floor; no second research round & B1 is two wire pads + two 100uF bulk caps with no protection by decision Q7; its loop/placement rules come from the verified B3 PVCC three-tier decoupling and dc-input hot-plug records; a 2nd round would cost \textasciitilde{}\$5-8 for little design delta \\
P3 & no TI Fig 37 PVCC RC damper (C58 100nF + R10 3.3R) added & C14 electrolytic ESR is the bulk damping per the verified dc-input-bulk-esr-damping-budget record; hot-plug ring margin 1.15x stays a P4-review watch item with RC snubber as fallback \\
P5 & U1 footprint fix: the 9 EasyEDA thermal vias (0.3 drill, 0.61 pad, 1.0 pitch 3x3) had empty pad numbers (no net -\textgreater{} 9 mask-bridge + 9 clearance errors); renumbered to pad 33 (GND) and set layers {*}.Cu only (no paste/mask on the vias: bottom tented, top inside the EP mask window) & thermal vias must be GND and join the bottom pour; record class-d-amp-powerpad-thermal-via-array asks component-side capping, which is impossible inside the EP window without plugging - via-in-pad plugging is the P9 alternative; 3x3 at 1.0 mm meets the \textgreater{}=6 count / \textless{}=1.2 mm pitch rule \\
P5 & wire-pad footprints: moved pin-1 silk corner 0.35 mm outside pad 1 (was over copper, 4 silk\_over\_copper warnings) & footprint-internal silk is a librarian fix; board\_init self-check now 0 violations \\
P4 & review W1/W2: swap C13 to a \textgreater{}=35 V low-ESR polymer/hybrid D6.3 part and C14 to a 35 V part with a stated ESR/ripple; ADD a unidirectional TVS (SMBJ15A-class, clamp \textless{} 30 V abs max) across J1 +12V/GND & reviewer: 26 V hot-plug ring exceeds C13's 25 V and C14's ESR cannot damp the bank; TVS bounds the ring below U1's 30 V abs max. Correctness over the reference design's part count (deviation for COMPARISON) \\
P4 & review W3: C1-C4 2.2 uF -\textgreater{} 1 uF X7R 50 V 0805, K tolerance (TI Fig 37 value; corner 2.7 Hz at 60k) & 2.2 uF x 60k (132 ms) with no source connected reaches the 420 ms DC-detect window at tolerance corners and fault-cycles on power-up; deviation from the reference design's 2.2 uF recorded \\
P4 & review W4: reverse polarity stays unprotected (Q7) but J1 gets a +12V/GND silk legend; W5: C11/C12 X5R 10 uF accepted & W4 accepted risk, legend is free; W5: \textasciitilde{}4-5 uF effective at 12 V is ample beside 200 uF bulk + 100 nF/1 nF HF caps \\
P4 & Schematic review W1-W3 fixed: C13 35V polymer C46550467, C14 FZ 35V C2992596 (stated ESR/ripple), C1-C4 1uF X7R C28323, D1 SMBJ15A C113988 across J1 & W1 26V ring \textgreater{} 25V C13 and C14 cannot damp it -\textgreater{} TVS clamps; W2 RVT had no ESR spec; W3 132ms input tau vs 420ms DC-detect -\textgreater{} 60ms \\
P6 & Accept 53.6 x 36.8 mm outline (2 layers) as a deviation from the 44 x 34 target & Thermal/decoupling layout first: 25x25 mm clear thermal copper around U1, PVCC stacks at the pins, 7.3 mm inductors need 9.6 mm column pitch and 3x6 mm speaker wire pads J3/J4 take their own column; closing \textasciitilde{}9 mm would need smaller inductors or wire pads, trading current rating/solderability for size \\
P6 & Strip body-outline silk from C0402/C0603/C0805/R0603 (lib + board) and widen IND courtyard to cover pads & 98 outline-vs-outline silk\_overlap warnings in the dense PVCC/bootstrap stacks would fail drc\_routed; passives are unpolarised, refdes kept; IND courtyard was smaller than its pads \\
P7 & P7-\textgreater{}P6: spread L3/L4 up and L1/L2 down \textasciitilde{}2.5 mm (grow outline height if needed) to fit the C15-C18 bootstrap caps in two columns and route OUT nets on F.Cu & OUT nets unroutable with current placement; caps-on-B.Cu alternative would cut the U1 heat-spreader pour edge with 8 vias; thermal first, area already a stated deviation \\
\bottomrule
\end{longtable}
}
\subsection*{Sheet plan}
\subsection*{PCB-0017-A stereo-class-d-amp - schematic sheet plan (P2)}
\subsubsection*{Ruling: ONE sheet, no hierarchy}
38 fitted parts around one IC, on one signal path (J2 -\textgreater{} U1 -\textgreater{} LC -\textgreater{} J3/J4) with one rail. A hierarchy would buy nothing but hierarchical pins and sheet-path net names; a single A3-class sheet shows the whole amplifier at a glance and the placement groups at P6 read straight off it. Root sheet = the only sheet: \texttt{PCB-0017-A\_stereo-class-d-amp.kicad\_sch} (same name as the KiCad project and the workspace directory, boards CLAUDE.md).

Net naming: power nets bare (\texttt{+12V}, \texttt{GND}); every other net is a local label and carries a leading \texttt{/} in the netlist (\texttt{/OUTPL}). These are the CANONICAL names: \texttt{constraints.json} uses them and the schematic must produce them exactly.

\subsubsection*{Drawing layout (left to right, mirroring the board)}
J2 and the four input caps at the far left; U1 in the middle drawn with pins 1-16 on its left and 17-32 on its right (the DAP footprint's own orientation: pin 1 top-left, PVCC 31/32 top-right, AVCC/PVCC 17-19 bottom-right); the configuration cluster (R1, R2, R3, C5, the MODSEL/AM/SYNC ties) below-left of U1; bootstrap and PVCC decoupling caps drawn at their pins on the right; the two LC filters to the right of U1 with the right channel above the left (matching pins 27/29 above 21/23); J4 then J3 at the far right; J1 with C13/C14 along the top right; TP1 on GND anywhere; one PWR\_FLAG each on \texttt{+12V} and \texttt{GND} (external rail, no on-board source - ERC needs them).

\subsubsection*{Pin-by-pin tie table (TI SLOS708G, DAP package - P3/P4 verify against the symbol)}
\begin{quote}\small\ttfamily\raggedright
\textbar{} pin \textbar{} name \textbar{} net \textbar{} tie / part \textbar{}\\
\textbar{}---\textbar{}---\textbar{}---\textbar{}---\textbar{}\\
\textbar{} 1 \textbar{} MODSEL \textbar{} `GND` \textbar{} BD mode \textbar{}\\
\textbar{} 2 \textbar{} SDZ \textbar{} `/SDZ` \textbar{} joined with FAULTZ; R2 100k to `+12V` \textbar{}\\
\textbar{} 3 \textbar{} FAULTZ \textbar{} `/SDZ` \textbar{} open-drain, same node as SDZ (auto-clearing fault) \textbar{}\\
\textbar{} 4 \textbar{} RINP \textbar{} `/RINP` \textbar{} C3 2.2 uF from `/IN\_R` \textbar{}\\
\textbar{} 5 \textbar{} RINN \textbar{} `/RINN` \textbar{} C4 2.2 uF to `GND` \textbar{}\\
\textbar{} 6 \textbar{} PLIMIT \textbar{} `/GVDD` \textbar{} direct tie, no limit \textbar{}\\
\textbar{} 7 \textbar{} GVDD \textbar{} `/GVDD` \textbar{} C5 1 uF to `GND` \textbar{}\\
\textbar{} 8 \textbar{} GAIN/SLV \textbar{} `/GAIN` \textbar{} R1 5.6k to `GND`, GVDD leg open: 20 dB master \textbar{}\\
\textbar{} 9 \textbar{} GND \textbar{} `GND` \textbar{} \textbar{}\\
\textbar{} 10 \textbar{} LINP \textbar{} `/LINP` \textbar{} C1 2.2 uF from `/IN\_L` \textbar{}\\
\textbar{} 11 \textbar{} LINN \textbar{} `/LINN` \textbar{} C2 2.2 uF to `GND` \textbar{}\\
\textbar{} 12 \textbar{} MUTE \textbar{} `/MUTE` \textbar{} R3 100k to `GND` (outputs enabled) \textbar{}\\
\textbar{} 13 \textbar{} AM2 \textbar{} `GND` \textbar{} 400 kHz (AM2/AM1/AM0 = 0/0/0) \textbar{}\\
\textbar{} 14 \textbar{} AM1 \textbar{} `GND` \textbar{} \textbar{}\\
\textbar{} 15 \textbar{} AM0 \textbar{} `GND` \textbar{} \textbar{}\\
\textbar{} 16 \textbar{} SYNC \textbar{} no connect \textbar{} output in master mode; ERC no-connect flag \textbar{}\\
\textbar{} 17 \textbar{} AVCC \textbar{} `+12V` \textbar{} C6 100 nF to `GND` at the pin \textbar{}\\
\textbar{} 18, 19 \textbar{} PVCC \textbar{} `+12V` \textbar{} C7 1 nF, C9 100 nF, C11 10 uF, C14 100 uF electrolytic \textbar{}\\
\textbar{} 20 \textbar{} BSNL \textbar{} `/BSNL` \textbar{} C15 220 nF to `/OUTNL` \textbar{}\\
\textbar{} 21 \textbar{} OUTNL \textbar{} `/OUTNL` \textbar{} L2 -\textgreater{} `/SPK\_L\_N` \textbar{}\\
\textbar{} 22 \textbar{} GND \textbar{} `GND` \textbar{} \textbar{}\\
\textbar{} 23 \textbar{} OUTPL \textbar{} `/OUTPL` \textbar{} L1 -\textgreater{} `/SPK\_L\_P` \textbar{}\\
\textbar{} 24 \textbar{} BSPL \textbar{} `/BSPL` \textbar{} C16 220 nF to `/OUTPL` \textbar{}\\
\textbar{} 25 \textbar{} GND \textbar{} `GND` \textbar{} \textbar{}\\
\textbar{} 26 \textbar{} BSNR \textbar{} `/BSNR` \textbar{} C17 220 nF to `/OUTNR` \textbar{}\\
\textbar{} 27 \textbar{} OUTNR \textbar{} `/OUTNR` \textbar{} L4 -\textgreater{} `/SPK\_R\_N` \textbar{}\\
\textbar{} 28 \textbar{} GND \textbar{} `GND` \textbar{} \textbar{}\\
\textbar{} 29 \textbar{} OUTPR \textbar{} `/OUTPR` \textbar{} L3 -\textgreater{} `/SPK\_R\_P` \textbar{}\\
\textbar{} 30 \textbar{} BSPR \textbar{} `/BSPR` \textbar{} C18 220 nF to `/OUTPR` \textbar{}\\
\textbar{} 31, 32 \textbar{} PVCC \textbar{} `+12V` \textbar{} C8 1 nF, C10 100 nF, C12 10 uF, C13 100 uF polymer \textbar{}\\
\textbar{} 33 \textbar{} PowerPAD \textbar{} `GND` \textbar{} via array to B.Cu \textbar{}
\end{quote}
\subsubsection*{Nets that leave a block (placement groups at P6, net classes at P5)}
\begin{quote}\small\ttfamily\raggedright
\textbar{} net \textbar{} from \textbar{} to \textbar{} class / note \textbar{}\\
\textbar{}---\textbar{}---\textbar{}---\textbar{}---\textbar{}\\
\textbar{} `+12V` \textbar{} J1.1 \textbar{} C13, C14, C11, C12, C9, C10, C7, C8, U1 17/18/19/31/32, R2 \textbar{} 4 A, \textgreater{}= 1.0 mm at 2 oz \textbar{}\\
\textbar{} `GND` \textbar{} J1.2, J2.2, TP1 \textbar{} everything; B.Cu pour + F.Cu pour \textbar{} 4 A return, heat spreader \textbar{}\\
\textbar{} `/IN\_L`, `/IN\_R` \textbar{} J2.3, J2.1 \textbar{} C1, C3 \textbar{} signal, keep \textgreater{}= 10 mm from L1-L4 \textbar{}\\
\textbar{} `/LINP`, `/RINP` \textbar{} C1, C3 \textbar{} U1 10, 4 \textbar{} 60 kOhm high-Z nodes, short \textbar{}\\
\textbar{} `/LINN`, `/RINN` \textbar{} C2, C4 \textbar{} U1 11, 5 \textbar{} high-Z, short \textbar{}\\
\textbar{} `/GVDD` \textbar{} U1 7 \textbar{} U1 6, C5 \textbar{} local \textbar{}\\
\textbar{} `/GAIN` \textbar{} U1 8 \textbar{} R1 \textbar{} local; sets gain and master mode, latched at start-up \textbar{}\\
\textbar{} `/SDZ` \textbar{} U1 2, 3 \textbar{} R2 \textbar{} local \textbar{}\\
\textbar{} `/MUTE` \textbar{} U1 12 \textbar{} R3 \textbar{} local \textbar{}\\
\textbar{} `/OUTPL`, `/OUTNL` \textbar{} U1 23, 21 \textbar{} L1, L2, C16, C15 \textbar{} 3 A switch nodes, \textgreater{}= 0.8 mm, symmetric pair, tight \textbar{}\\
\textbar{} `/OUTPR`, `/OUTNR` \textbar{} U1 29, 27 \textbar{} L3, L4, C18, C17 \textbar{} same \textbar{}\\
\textbar{} `/BSPL`, `/BSNL`, `/BSPR`, `/BSNR` \textbar{} U1 24, 20, 30, 26 \textbar{} C16, C15, C18, C17 \textbar{} bootstrap, loop local to its pin pair \textbar{}\\
\textbar{} `/SPK\_L\_P`, `/SPK\_L\_N` \textbar{} L1, L2 \textbar{} C19, C20, R4, R5, J3.1, J3.2 \textbar{} 3 A, \textgreater{}= 0.8 mm \textbar{}\\
\textbar{} `/SPK\_R\_P`, `/SPK\_R\_N` \textbar{} L3, L4 \textbar{} C21, C22, R6, R7, J4.1, J4.2 \textbar{} 3 A \textbar{}\\
\textbar{} `/ZL\_P`, `/ZL\_N`, `/ZR\_P`, `/ZR\_N` \textbar{} R4, R5, R6, R7 \textbar{} C23, C24, C25, C26 \textbar{} damping-branch midpoints, local \textbar{}
\end{quote}
\subsubsection*{Refdes ranges (one sheet; fixed so a later renumber lands in the same places)}
\begin{quote}\small\ttfamily\raggedright
\textbar{} class \textbar{} range \textbar{} assigned \textbar{}\\
\textbar{}---\textbar{}---\textbar{}---\textbar{}\\
\textbar{} U \textbar{} U1-U9 \textbar{} U1 TPA3118D2DAPR \textbar{}\\
\textbar{} L \textbar{} L1-L9 \textbar{} L1 OUTPL, L2 OUTNL, L3 OUTPR, L4 OUTNR - 10 uH 7.3 x 6.6 mm \textbar{}\\
\textbar{} R \textbar{} R1-R19 \textbar{} R1 5.6k gain; R2 100k SDZ pull-up; R3 100k MUTE pull-down; R4 SPK\_L\_P, R5 SPK\_L\_N, R6 SPK\_R\_P, R7 SPK\_R\_N damping 3.3R 0603 \textbar{}\\
\textbar{} C \textbar{} C1-C39 \textbar{} C1 LINP / C2 LINN / C3 RINP / C4 RINN 2.2 uF 0805; C5 1 uF GVDD 0603; C6 100 nF AVCC 0603; C7, C8 1 nF C0G 0603 (PVCC 18/19, 31/32); C9, C10 100 nF 0603; C11, C12 10 uF 0805; C13 100 uF polymer (31/32), C14 100 uF electrolytic 35 V (18/19); C15 BSNL, C16 BSPL, C17 BSNR, C18 BSPR 220 nF 0805; C19 SPK\_L\_P, C20 SPK\_L\_N, C21 SPK\_R\_P, C22 SPK\_R\_N 680 nF 0805; C23-C26 10 nF 0402 damping (same order) \textbar{}\\
\textbar{} J \textbar{} J1-J9 \textbar{} J1 power 2 pads; J2 line in 3 pads (R, GND, L); J3 left spk 2 pads; J4 right spk 2 pads - all SMD wire pads, no connector body \textbar{}\\
\textbar{} TP \textbar{} TP1-TP9 \textbar{} TP1 GND \textbar{}\\
\textbar{} `\#PWR` \textbar{} `pwr\_base = 100` -\textgreater{} `\#PWR0101...` \textbar{} `+12V`, `GND` symbols \textbar{}\\
\textbar{} `\#FLG` \textbar{} `\#FLG0101...` \textbar{} one PWR\_FLAG on `+12V`, one on `GND` \textbar{}
\end{quote}
\subsubsection*{Values and ratings the schematic must carry}
\begin{quote}\small\ttfamily\raggedright
\textbar{} ref \textbar{} value \textbar{} rating / dielectric \textbar{} why \textbar{}\\
\textbar{}---\textbar{}---\textbar{}---\textbar{}---\textbar{}\\
\textbar{} C1-C4 \textbar{} 2.2 uF \textbar{} X7R, 50 V (25 V min), 0805 \textbar{} the reference design's value; 1.2 Hz corner at 60 k \textbar{}\\
\textbar{} C5 \textbar{} 1 uF \textbar{} X7R 25 V 0603 \textbar{} GVDD LDO decoupling, datasheet-mandated \textbar{}\\
\textbar{} C6 \textbar{} 100 nF \textbar{} X7R 50 V 0603 \textbar{} AVCC \textbar{}\\
\textbar{} C7, C8 \textbar{} 1 nF \textbar{} C0G/NP0 50 V 0603 \textbar{} switching-edge loop \textbar{}\\
\textbar{} C9, C10 \textbar{} 100 nF \textbar{} X7R 50 V 0603 \textbar{} mid-frequency \textbar{}\\
\textbar{} C11, C12 \textbar{} 10 uF \textbar{} X7R 25 V 0805 \textbar{} P1 decision (Basic 0805, not 1206) \textbar{}\\
\textbar{} C13 \textbar{} 100 uF \textbar{} Al-polymer 25 V, ESR \textless{}= 30 mOhm, \textgreater{}= 2 A rms \textbar{} ripple \textbar{}\\
\textbar{} C14 \textbar{} 100 uF \textbar{} Al-electrolytic 35 V low-ESR \textbar{} ripple + damps hot-plug ring \textbar{}\\
\textbar{} C15-C18 \textbar{} 220 nF \textbar{} X7R 50 V 0805 \textbar{} bootstrap, one per output \textbar{}\\
\textbar{} C19-C22 \textbar{} 680 nF \textbar{} X7R 50 V 0805 \textbar{} LC filter, f0 61 kHz with 10 uH \textbar{}\\
\textbar{} C23-C26 \textbar{} 10 nF \textbar{} X7R 50 V 0402 \textbar{} damping branch \textbar{}\\
\textbar{} L1-L4 \textbar{} 10 uH \textbar{} Isat \textgreater{}= 4 A (5.5 A lead), Irms \textgreater{}= 3 A, DCR \textless{}= 70 mOhm, 7.3 x 6.6 mm shielded \textbar{} 3.5 A peak, 7.5 A OC trip \textbar{}\\
\textbar{} R1 \textbar{} 5.6k \textbar{} 1 \%, 0603 \textbar{} 20 dB / master; latched at power-up \textbar{}\\
\textbar{} R2, R3 \textbar{} 100k \textbar{} 1 \%, 0603 \textbar{} SDZ+FAULTZ pull-up, MUTE pull-down \textbar{}\\
\textbar{} R4-R7 \textbar{} 3.3R \textbar{} 1 \%, 0603 (\textgreater{}= 100 mW) \textbar{} damping \textbar{}
\end{quote}
Full architecture narrative (not inlined; see the artifact index): \texttt{architecture/blocks.md}, \texttt{architecture/power\_tree.md}, \texttt{architecture/stackup.md}.
\section{Schematic}
Gate \texttt{erc} (phase P4): \textbf{pass} after 2 attempt(s); last 2026-09-27T03:45:59, failing 0 of 0.
The full schematic PDF follows.
\includepdf[pages=-]{reports/schematic.pdf}
\section{Layout}
Gate \texttt{place} (phase P6): \textbf{pass} after 31 attempt(s); last 2026-09-27T08:55:02, failing 0 of 0.

Gate \texttt{drc\_routed} (phase P7): \textbf{pass} after 14 attempt(s); last 2026-09-27T08:55:04, failing 0 of 0.
\begin{center}
\includegraphics[width=0.72\textwidth]{reports/PCB-0017-A_stereo-class-d-amp_top.png}\\
{\small\texttt{reports/PCB-0017-A\_stereo-class-d-amp\_top.png}}
\end{center}
\begin{center}
\includegraphics[width=0.72\textwidth]{reports/PCB-0017-A_stereo-class-d-amp_bottom.png}\\
{\small\texttt{reports/PCB-0017-A\_stereo-class-d-amp\_bottom.png}}
\end{center}
\section{Verification}
\emph{Pending --- produced at P8.}
\section{DFM and Fabrication}
\emph{Pending --- produced at P9.}
\section{Run Record (Appendix)}
\subsection*{Phase digests}
\paragraph*{P0}
\begin{quote}\small\ttfamily\raggedright
\# P0 digest\\
- requirements.md written by requirements-analyst (fable), check\_requirements pass, no mode token.\\
- 12 open questions answered by the orchestrator as delegate (defaults accepted); Q6 current 3 A PROVISIONAL.\\
- Size soft target \textless{}= 44 x 34 mm, 2 layers; thermal pad / heat-sinking outranks size.\\
- Assembly ASSUMED JLC PCBA top side; qty 5; bare wire pads J1-J4 + TP1; no protection, no holes.
\end{quote}
\paragraph*{P1}
\begin{quote}\small\ttfamily\raggedright
\# P1 digest\\
- Roster: component-scout (sonnet), reference-design (sonnet, incl. both reference-design sources), power-architect (opus).\\
- DEVIATION: TPA3116D2 is pad-UP only (DAD); swapped to pin-compatible pad-down TPA3118D2DAPR (C46497) so the board copper sinks it.\\
- +12V 4 A (3.6 A at 10\% THD, 4 ohm); outputs 3 A; chip loss \textasciitilde{}3 W (4 ohm 1\%), \textasciitilde{}1.2 W (8 ohm), \textasciitilde{}0.8 W music.\\
- 44x34 mm holds 8 ohm full power and 4 ohm music; sustained 4 ohm full sine needs \textasciitilde{}50x50 mm (not required).\\
- L = SLO0630H100MTT 7.3x6.6 mm (Isat 5.5 A) - the size driver; 680 nF filter caps; 3.3R+10nF damping; 2x100 uF (polymer+electrolytic).\\
- 2 oz copper preferred; PLIMIT-\textgreater{}GVDD, FAULTZ-\textgreater{}SDZ, 20 dB gain (5.6k), 100k pull-ups.\\
- Quilter source had no board-specific content; JLC source: 44x34, 2L, 45 fp / 40 fitted, 0 ERC/DRC, never fabricated.
\end{quote}
\paragraph*{P2}
\begin{quote}\small\ttfamily\raggedright
\# P2 digest\\
- Architect (fable/high): 4 blocks - B1 12 V/4 A wire-pad input + 100 uF polymer \& 100 uF electrolytic bulk; B2 stereo line-in, 2.2 uF X7R coupling; B3 TPA3118D2DAPR BTL, BD, 400 kHz, 20 dB, master; B4 10 uH + 680 nF LC per leg, 3.3R+10nF damping.\\
- One schematic sheet, no hierarchy (38 fitted parts); canonical net names in sheets.md, used by constraints.json.\\
- Stackup: 2 layers, 2 oz outer preferred; B.Cu one GND pour with the PowerPAD thermal-via array into it.\\
- Size driver: four 7.3x6.6 mm inductors; target \textless{}= 44x34 mm, placement earns the outline (provisional, then fit).\\
- constraints\_lint exit 0 (11 unitless operating-point warnings, informational).\\
- Coverage P2: 4/4 slots gap (no checklists for these topologies) -\textgreater{} 4 research tasks opened automatically.
\end{quote}
\paragraph*{P3}
\emph{(no digest recorded)}
\paragraph*{P4}
\emph{(no digest recorded)}
\paragraph*{P5}
\emph{(no digest recorded)}
\paragraph*{P6}
\emph{(no digest recorded)}
\paragraph*{P7}
\emph{(no digest recorded)}
\subsection*{Gate attempt history}
{\small
\begin{longtable}{llll}
\toprule
\textbf{Timestamp} & \textbf{Gate} & \textbf{Status} & \textbf{Failing/Total} \\
\midrule
\endhead
2026-09-27T03:19:32 & erc & pass & 0/0 \\
2026-09-27T03:45:59 & erc & pass & 0/0 \\
2026-09-27T04:10:46 & place & fail & 49/49 \\
2026-09-27T04:11:28 & place & fail & 49/49 \\
2026-09-27T04:11:48 & place & fail & 49/49 \\
2026-09-27T04:13:04 & place & fail & 21/21 \\
2026-09-27T04:14:33 & place & fail & 3/3 \\
2026-09-27T04:15:31 & place & fail & 1/1 \\
2026-09-27T04:15:55 & place & fail & 1/1 \\
2026-09-27T04:16:16 & place & pass & 0/0 \\
2026-09-27T04:17:57 & place & pass & 0/0 \\
2026-09-27T04:19:59 & place & pass & 0/0 \\
2026-09-27T04:52:12 & drc\_routed & fail & 56/56 \\
2026-09-27T05:29:04 & place & fail & 5/5 \\
2026-09-27T05:47:45 & place & fail & 2/2 \\
2026-09-27T05:51:46 & place & fail & 1/1 \\
2026-09-27T05:52:46 & place & fail & 1/1 \\
2026-09-27T05:53:10 & place & pass & 0/0 \\
2026-09-27T05:54:01 & place & pass & 0/0 \\
2026-09-27T05:54:06 & drc\_routed & fail & 7/7 \\
2026-09-27T05:54:43 & place & pass & 0/0 \\
2026-09-27T05:54:48 & drc\_routed & fail & 6/6 \\
2026-09-27T05:55:08 & place & pass & 0/0 \\
2026-09-27T06:01:34 & place & fail & 5/5 \\
2026-09-27T06:02:16 & place & pass & 0/0 \\
2026-09-27T06:02:17 & place & pass & 0/0 \\
2026-09-27T06:02:23 & drc\_routed & fail & 6/6 \\
2026-09-27T06:23:13 & place & pass & 0/0 \\
2026-09-27T06:23:17 & drc\_routed & fail & 3/3 \\
2026-09-27T06:24:44 & drc\_routed & fail & 2/2 \\
2026-09-27T06:25:19 & drc\_routed & pass & 0/0 \\
2026-09-27T06:25:22 & place & pass & 0/0 \\
2026-09-27T06:34:02 & place & pass & 0/0 \\
2026-09-27T06:34:07 & drc\_routed & pass & 0/0 \\
2026-09-27T06:35:39 & dfm & pass & 0/0 \\
2026-09-27T06:37:18 & place & pass & 0/0 \\
2026-09-27T06:37:23 & drc\_routed & pass & 0/0 \\
2026-09-27T06:38:29 & dfm & pass & 0/0 \\
2026-09-27T06:55:35 & verify & fail & 8/31 \\
2026-09-27T06:57:55 & verify & fail & 8/31 \\
2026-09-27T07:21:41 & drc\_routed & pass & 0/0 \\
2026-09-27T07:21:43 & place & pass & 0/0 \\
2026-09-27T07:26:10 & dfm & pass & 0/0 \\
2026-09-27T07:26:25 & verify & fail & 8/31 \\
2026-09-27T07:49:27 & place & pass & 0/0 \\
2026-09-27T07:50:41 & drc\_routed & pass & 0/0 \\
2026-09-27T07:54:33 & dfm & pass & 0/0 \\
2026-09-27T08:07:10 & place & pass & 0/0 \\
2026-09-27T08:07:14 & drc\_routed & pass & 0/0 \\
2026-09-27T08:13:09 & dfm & pass & 0/0 \\
2026-09-27T08:13:19 & verify & fail & 8/31 \\
2026-09-27T08:51:08 & place & fail & 7/7 \\
2026-09-27T08:51:11 & drc\_routed & pass & 0/0 \\
2026-09-27T08:51:21 & dfm & pass & 0/0 \\
2026-09-27T08:54:10 & place & pass & 0/0 \\
2026-09-27T08:55:02 & place & pass & 0/0 \\
2026-09-27T08:55:04 & drc\_routed & pass & 0/0 \\
2026-09-27T08:55:11 & dfm & pass & 0/0 \\
2026-09-27T08:57:46 & dfm & pass & 0/0 \\
2026-09-27T08:57:51 & verify & fail & 8/32 \\
\bottomrule
\end{longtable}
}
\subsection*{Issues}
\emph{(no entries)}
\subsection*{Human checkpoints}
{\small
\begin{longtable}{lllp{2.6cm}p{6.4cm}}
\toprule
\textbf{Id} & \textbf{Phase} & \textbf{Status} & \textbf{When} & \textbf{Note} \\
\midrule
\endhead
1 & P2 & approved & 2026-09-27T03:07:31 & delegate (owner: decide for yourself). Coverage: C46497 covered; B2/B3/B4 provisional, all 4 research tasks closed verified; B1 gap accepted by decision. TPA3118D2 DAP pad-down instead of 3116 DAD pad-up (needs heatsink); 2L 2oz; 4 A input; riskiest: 2L thermal at 4 ohm sustained, hot-plug ring margin 1.15x. \\
2 & P4 & approved & 2026-09-27T03:46:57 & delegate (owner: decide for yourself). Schematic 44 parts/29 nets, ERC 0/0, netlist\_audit pass; review 0 err/5 warn, W1-W3 fixed, W4 silk legend at layout, W5 X5R accepted. Deviation from the reference design: TPA3118D2 pad-down + SMBJ15A TVS added + 2x100uF bulk. \\
3 & P6 & approved & 2026-09-27T04:35:06 & delegate (owner: decide for yourself). Place gate PASS, DRC 0 after silk fix, 53.6x36.8 mm 2L accepted as deviation (decision). \\
4 & P8 & not reached &  &  \\
5 & P10 & not reached &  &  \\
\bottomrule
\end{longtable}
}
\subsection*{Generation cost}
Model cost of generating this board: \textbf{\$97.57}.

{\small
\begin{longtable}{p{10cm}r}
\toprule
\textbf{Step} & \textbf{Cost} \\
\midrule
\endhead
Before the workspace (brief and skill read) & \$0.50 \\
Brief and intake & \$1.16 \\
Research & \$3.84 \\
Architecture & \$14.22 \\
Schematic & \$3.10 \\
Placement & \$4.88 \\
Routing & \$0.42 \\
Verification & \$0.51 \\
DFM, order files and guide & \$0.11 \\
Not in the transcript (side calls) & \$68.83 \\
\bottomrule
\end{longtable}
}
\section{Artifact Index}
{\small
\begin{longtable}{p{9cm}rp{4cm}}
\toprule
\textbf{File} & \textbf{Size} & \textbf{Note} \\
\midrule
\endhead
\texttt{state.json} & 69,768 B &  \\
\texttt{kicad/PCB-0017-A\_stereo-class-d-amp.kicad\_pcb} & 261,364 B &  \\
\texttt{kicad/PCB-0017-A\_stereo-class-d-amp.kicad\_sch} & 151,584 B &  \\
\texttt{fab/PCB-0017-A\_stereo-class-d-amp\_gerbers.zip} & 45,371 B &  \\
\texttt{fab/BOM.csv} & 743 B &  \\
\texttt{fab/CPL.csv} & 1,334 B &  \\
\texttt{brief/brief.md} & 953 B &  \\
\texttt{requirements.md} & 9,052 B &  \\
\texttt{architecture/blocks.md} & 8,474 B &  \\
\texttt{architecture/constraints.json} & 15,472 B &  \\
\texttt{architecture/power\_tree.md} & 6,460 B &  \\
\texttt{architecture/sheets.md} & 6,644 B &  \\
\texttt{architecture/stackup.md} & 4,159 B &  \\
\texttt{reports/PCB-0017-A\_stereo-class-d-amp\_bottom.png} & 38,076 B &  \\
\texttt{reports/PCB-0017-A\_stereo-class-d-amp\_top.png} & 230,278 B &  \\
\texttt{reports/board\_init.json} & 1,038 B &  \\
\texttt{reports/bom\_cpl.json} & 22,913 B &  \\
\texttt{reports/check\_thermal.json} & 2,013 B &  \\
\texttt{reports/cost.json} & 6,761 B &  \\
\texttt{reports/fab\_export.json} & 2,359 B &  \\
\texttt{reports/gate-dfm.json} & 1,552 B &  \\
\texttt{reports/gate-drc\_routed.json} & 1,022 B &  \\
\texttt{reports/gate-erc.json} & 822 B &  \\
\texttt{reports/gate-place-optionA.json} & 4,103 B &  \\
\texttt{reports/gate-place-optionA3.json} & 4,171 B &  \\
\texttt{reports/gate-place.json} & 1,407 B &  \\
\texttt{reports/gate-verify.json} & 8,573 B &  \\
\texttt{reports/guide\_facts.json} & 12,627 B &  \\
\texttt{reports/placement-notes.md} & 2,958 B &  \\
\texttt{reports/review-schematic.json} & 6,263 B &  \\
\texttt{reports/review-schematic.md} & 14,056 B &  \\
\texttt{reports/review-verify.md} & 9,422 B &  \\
\texttt{reports/schematic.pdf} & 149,082 B &  \\
\texttt{reports/top.net} & 73,871 B &  \\
\texttt{reports/verify-waivers.proposed.json} & 6,139 B &  \\
\bottomrule
\end{longtable}
}
\end{document}
